\begin{textsample}{POS Dim 3 – 2009\_c5 – Score 1.00 – 2009\_c5/enade\_2009\_est\_f2e70bb570.txt}  \label{ex:f3_pos_139}
Capítulo 5 Distribuição dos Conceitos 5.1 Panorama nacional da distribuição dos conceitos A Tabela 5.1 apresenta a quantidade e o percentual de cursos de ESTATÍSTICA participantes do ENADE/2009 que se situaram em cada uma das faixas de conceitos, distribuídos de acordo com a Região em que se \textbf{localizam} os cursos das instituições de educação superior, mantenedoras dos cursos responsáveis pela oferta.
O Gráfico 5.1 apresenta a mesma informação sob a forma de diagrama de caixa para uma melhor visualização da dispersão dos conceitos obtidos.
No diagrama de caixa, a caixa propriamente dita fica limitada pelos quartis.
A barra dentro da caixa em negrito representa a mediana.
Linhas verticais unem os quartis aos valores extremos.
Se estes valores extremos se configurarem como valores atípicos, estes aparecem como pontos isolados e as linhas terminam no maior ( ou menor ) valor ainda não considerado atípico.
Eventuais valores atípicos aparecem como pequenos círculos com a Categoria Administrativa correspondente.
Valores considerados muito atípicos foram representados por estrelas.
Um diagrama de caixa para poucos dados ou com dados muito concentrados pode ter alguns pontos coincidentes.
Por exemplo, se mais de metade dos cursos de uma Região receberem conceito 5, o valor extremo superior, o quartil superior e a mediana coincidirão.
Se o conjunto só contém um dado, este apareceria denotado como a barra da mediana.
Se o conjunto apresenta todas as notas iguais, a notação é a mesma, somente uma barra horizontal.
Observando -se os dados da Tabela 5.1, percebe -se que, em termos nacionais, predominaram os cursos que receberam conceitos 2 e 3, correspondentes, respectivamente a 20,0 \% e a 32,0 \% do total nacional.
Foram classificados com conceito mínimo 16,0 \% dos cursos.
Conceitos 4 e 5 foram alocados a, respectivamente, 4,0 \% e 8,0 \%.
Dos cursos, apenas 5 ficaram sem conceito, correspondendo a 20,0 \% do total.
Considerando -se separadamente as cinco Regiões observou -se que, da Região Norte, 2 cursos participaram, dos quais 1 obteve conceito 2 e o outro, conceito 3.
Na Região Nordeste, 7 cursos participaram, dos quais, 4 obtiveram conceito mínimo, 1 o conceito e 1 o conceito 3.
Nenhum curso desta Região obteve conceito máximo e 1 ficou sem conceito.
Na Região Sudeste, 11 cursos participaram, dos quais, 4 obtiveram conceito e 2 o conceito 2.
O conceito máximo foi alcançado por 2 cursos, nenhum obteve conceito mínimo e 3 cursos ficaram sem conceito.
Na Região Sul, dos 3 cursos avaliados, 2 classificaram -se com conceito 3 e 1 com conceito 2.
Na Região Centro-Oeste, 2 cursos participaram, 1 teve o conceito 4 e 1 curso não obteve conceito. 5.2 Conceitos por Categoria Administrativa e por Região A Tabela 5.2 apresenta a distribuição dos cursos participantes do ENADE/2009 de ESTATÍSTICA, por Categoria Administrativa, de acordo com os conceitos por eles alcançados, segundo as Regiões.
O Gráfico 5.2 apresenta a mesma informação sob a forma de diagrama de caixa para uma melhor visualização da dispersão dos conceitos obtidos.
Somente eventuais valores atípicos aparecem com o identificador da Região na qual o curso foi habilitado.
De acordo com as informações apresentadas, a grande maioria, 76,0 \%, é de cursos em instituições federais, 16,0 \% dos cursos eram vinculados a instituições estaduais e 8,0 \% a instituições privadas.
Apenas instituições federais apresentaram conceito 5.
O conceito modal para esta categoria foi 3, com 31,6 \% dos cursos, 2 cursos em instituição federal receberam o conceito mínimo e 5 ficaram sem conceito.
Entre os 4 cursos de instituições estaduais, 2 receberam conceito 3, 1 o conceito 2 e 1 o conceito mínimo.
Os 2 cursos nas instituições privadas tiveram conceitos 1 e 2.
Na Região Norte, os 2 cursos eram de instituições federais e receberam conceitos 2 e 3.
Na Região Nordeste, dos 7 cursos 5 eram de instituições federais, 1 de estadual e 1 de privada.
Os cursos em instituição federal receberam conceito 1 ( 2 cursos ) e conceitos 2 e 3 ( 1 cursos cada ).
Ficou sem conceito 1 cursos em instituição federal.
Receberam conceito mínimo o único curso de instituição estadual e o único de instituição privada.
Na Região Sudeste, dos 8 cursos em instituição federal, 3 ficaram com conceito 3, 2 com conceito máximo e 3 ficaram sem conceito.
Os 2 cursos em instituição estadual ficaram com conceito 2 e 3.
O único cursos em instituição privada obteve conceito 2.
Na Região Sul, dos 2 cursos em instituições federais, 1 ficou com conceito 3 e o outro com conceito 2.
O único curso em instituição estadual ficou com conceito 3.
Na Região Centro-Oeste, os 2 cursos participantes estavam vinculados a instituições federais, 1 obteve conceito 4 e outro ficou sem conceito. 5.3 Conceitos por Organização Acadêmica e por Região Na Tabela 5.3 encontra -se a distribuição dos conceitos atribuídos aos cursos participantes do ENADE/2009 de ESTATÍSTICA, por Organização Acadêmica, segundo as Regiões.
O Gráfico 5.3 apresenta a mesma informação sob a forma de diagrama de caixa para uma melhor visualização da dispersão dos conceitos obtidos.
Somente eventuais valores atípicos apareceram com o identificador da Categoria Administrativa.
Os cursos desta área estão principalmente alocados em Universidades ( 88,0 \% ), depois em Faculdades ( 8,0 \% ) e Centros Universitários ( 4,0 \% ).
Não existe presença de Centros Federais de Educação Tecnológica ( CEFET ) e Institutos Federais de Educação, Ciência e Tecnologia ( IFECT ) nesta área.
Nas Universidades, o conceito modal foi 3 ( 31,8 \% ), 3 cursos obtiveram conceito mínimo e 2 cursos o conceito máximo, sendo que 5 cursos ficaram sem conceito.
Nos Centros Universitários o único curso participante obteve conceito 2.
Nas Faculdades, os 2 cursos participantes obtiveram conceitos 1 e 3.
Considerando -se separadamente as cinco Regiões, observa -se que, a Região Norte apresentou 2 cursos em Universidades, com conceitos 2 e 3.
A Região Nordeste apresentou 6 cursos em Universidades e 1 em Faculdade.
Dos 6 cursos em Universidades, 3 obtiveram o conceito mínimo, 1 o conceito 2, 1 o conceito 3 e 1 ficou sem conceito.
O único curso em Faculdade obteve conceito mínimo.
A Região Sudeste apresentou 9 cursos em Universidades, 1 em Centro Universitário e 1 em Faculdade.
Dos 9 cursos em Universidade, 1 obteve conceito 2, 3 o conceito 3, 2 conceito máximo e 3 ficaram sem conceito.
Os únicos cursos em Centros Universitários e Faculdades obtiveram conceitos 2 e 3, respectivamente.
A Região Sul apresentou os 3 cursos em Universidades, sendo que 1 obteve o conceito 2 e 2 obtiveram conceito 3.
Assim como a Região Sul, a Região Centro-Oeste apresentou os 2 cursos em Universidades, sendo que 1 obteve o conceito 4 e o outro ficou sem conceito.

% matched lemmas: localizar
\end{textsample}
